% Draft 2026-11-03: official setup. \todo{} values come from results/official/analysis.json. The pilot's text is at tag pilot-v1.
\section{Results}
\label{sec:results}

\subsection{RQ1: Does generated memory beat no memory?}
\IfFileExists{sections/generated-results.tex}{% Generated by scripts/results_table.py from results/official/analysis.json. Do not edit by hand.
\begin{table}[t]
\centering
\caption{Strict passes per task: no memory (A0, 3 attempts) and Hermes memory (HM, 3 runs of 3 attempts).}
\label{tab:main}
\footnotesize
\begin{tabular}{@{}lcc@{}}
\toprule
Task & A0 & HM \\
\midrule
APEX curtailment & 0/3 & 0/9 \\
APEX market share & 0/3 & 3/9 \\
APEX seat surplus & 0/3 & 1/9 \\
AutoBench contract & 1/3 & 0/9 \\
AutoBench marketing & 2/3 & 4/9 \\
AutoBench proposal & 0/3 & 0/9 \\
DAPT & 0/3 & 0/9 \\
EntOps CSM & 2/3 & 4/9 \\
EntOps teams & 0/3 & 0/9 \\
Word2Vec & 0/3 & 5/9 \\
awilix explicit & 0/3 & 0/9 \\
awilix implicit & 0/3 & 0/9 \\
cattrs & 0/3 & 2/9 \\
cookie store & 1/3 & 0/9 \\
customer delivery & 0/3 & 0/9 \\
\midrule
Pooled & 6/45 & 19/135 \\
\bottomrule
\end{tabular}
\end{table}

\newcommand{\resRuns}{3}
\newcommand{\resMeanDiff}{+0.007}
\newcommand{\resCI}{[-0.12, +0.14]}
\newcommand{\resP}{0.500}
\newcommand{\resUp}{4}
\newcommand{\resDown}{4}
\newcommand{\resSensNoRefP}{0.344}
\newcommand{\resSensUnavailP}{0.500}

Table~\ref{tab:main} gives strict passes per task.
Memory changes the per-task pass rate by \resMeanDiff{} on average (exact 95\% interval \resCI), with \resUp{} tasks up and \resDown{} down; the pre-specified one-sided test gives $p=\resP$.
Without the 4 tasks whose reference solution scores 0, $p=\resSensNoRefP$; no counted trial gave up on the model server, so the pre-specified check that excludes such trials gives the same result.
Three of four memory runs have trials: run 2's generation hit the time limit (Table~\ref{tab:chain}), and run 4 replaced it by amendment; without run 4, the mean is $-0.01$ and $p=0.63$; without Word2Vec it is $-0.03$ ($-0.14$ to $+0.08$, $p=0.77$, post hoc).

Runs vary: 7, 4 and 8 of 45 passed, and EntOps CSM passed 2, 0 and 2 of 3. Only Word2Vec improved in every run (0 of 3 without memory; 2, 2 and 1 of 3 in runs 1, 3 and 4). 

On the pre-specified secondary, the share of hidden tests passed, memory changed the per-task rate by $+0.06$ ($-0.07$ to $+0.20$, $p=0.25$); Word2Vec rose from 0\% to 72--94\% of its tests with memory. In an exploratory placebo arm, run 1's memory with each file's text replaced by unrelated prose of the same length, Word2Vec passed 0 of 3 and the arm was 0.04 below no memory ($p=0.89$). The Word2Vec gain did not appear with the placebo (one task, 3 attempts).}{}

\subsection{RQ2: Where can memory's value be lost?}
Table~\ref{tab:chain} summarizes each link, then two problems that cut across links, serving and build. Unless marked, a finding comes from the public code, documentation or task files, or from running them, and does not depend on our model.

\begin{table}[t]
\centering
\caption{Where memory's value can be lost. S: source or documentation. O: observed by running the official pipeline. L: only for our local model.}
\label{tab:chain}
\footnotesize
\begin{tabular}{@{}p{0.17\columnwidth}p{0.65\columnwidth}c@{}}
\toprule
Link & Finding & Basis \\
\midrule
Corpus & Runs with prompts resembling a target task are very uneven: none for 8 of 15 tasks, 35 for DAPT with none passing; v1.1 does not change this. Same-repository runs with different prompts are not counted. & O \\
Generation & Memory is exported only when generation ends, so a timeout leaves none. & S, O \\
Generation & Three generations took 21 to 27 minutes of the 6 hours allowed. Run 2's generation delegated work to sub-agents that queued on the single model slot (median wait 10 minutes) and hit the limit with no memory. On a second GPU type that decodes about 1.5$\times$ slower, with identical serving settings, a generation timed out the same way, so the 6-hour limit may tie memory to serving capacity. One generation printed about 0.2\% of the corpus to itself (290\,KB of 126\,MB). & L \\
Harness & Index caps are fixed regardless of the context window: 2{,}500 tokens in Codex and 200 lines or 25\,KB in Claude Code, so 21\% and 48\% of the same entry file is visible at start. Hermes shows the whole entry file and a skill index, but that was 1 of 13 files of our run-1 memory; the rest must be opened. & S, O \\
Harness & At the highest effort, Hermes cancels a model call after 300\,s, but the cancel does not stop the server, so retries queue behind abandoned calls; after 5 it gives up, and the trial scores 0 and counts as valid. Hermes's post-answer memory review is cut when the trial's container stops. Text configured as an extra system prompt reaches every call but is never saved to the session log, so checks of what the model saw that read saved logs under-report it. & S, O \\
Agent & In run 1 the agent opened a memory file in 15 of 45 trials, never on 7 of 15 tasks. Of its 7 passes, 2 followed an open (both Word2Vec); 5 passed without opening memory beyond what was shown. In an exploratory arm, one added sentence telling the agent to read its memory raised opens to 27 of 45 trials and moved no score (8 vs 7 passes of 45); showing every file at the start, or none, gave 7 and 9 passes of 45. & L \\
Recency & No released task needs newer corpus evidence to be solved, so whether memory handles newer evidence cannot be tested on the released tasks. & O \\
Grader & The reference solution scores 0 on 4 of 15 tasks, as the maintainers document: on cattrs and cookie store it passes the functional tests but not the hidden preference checks, and on both EnterpriseOps tasks its trace is empty. Without memory, our agent met every hidden preference on cookie store in 1 of 3 counted attempts and passed EnterpriseOps CSM in 2 of 3; on cattrs it did so once in the discarded first control. These preferences can therefore be met without memory; with run-1 memory, no cookie-store attempt met all of them. EnterpriseOps Teams passed exactly 3 of 4 checks in all 6 attempts of the control and run 1. & O \\
Grader & In the grading code, a failed LLM judge yields reward 0 and a valid trial, indistinguishable from an agent failing. The task files and the CLI name different judges, and the leaderboard does not say which was used. & S \\
Serving & Harbor's egress proxy, which enforces no-network tasks, drops a model response not started within 15\,s. Any slower model hits it; the agent gives up and the trial counts as a valid failure (3 of 41 trials in a discarded first control). A malformed tool call from the model surfaces as an HTTP 500 (1 in about 4{,}100 calls, recovered by retry). & O \\
Build & Harbor rebuilds images per trial, and 12 of 15 tasks have install commands in their builds. One task broke without any change on the benchmark's side when a dependency released on Nov 1; another built to a different image from identical files. & S, O \\
\bottomrule
\end{tabular}
\end{table}

\subsection{RQ3: What can a reader conclude from the leaderboard?}
The leaderboard~\cite{frontiermemory} reports a score and a task cost for each entry. Its footnote says each score is a mean $\pm$1 sample standard deviation across 3 independent runs, but the rows show no standard deviations, and the page gives no no-memory row, per-task results, latency, tokens, memory size or judge. The rules do refer to no-memory rows for comparison.
On Nov 11, 2026, the 12 scores in the model table were all multiples of 1/45, which fits 15 tasks in each of 3 runs, and the top five were 33, 33, 32, 32 and 31 of 45. Treat each score as 45 independent attempts at its published rate. That is an upper bound on noise, because the same 15 tasks recur in every run, and it ignores that the leaderboard takes the higher of two task models. Under it, a rerun would reverse each adjacent pair among the top four with probability 0.41 to 0.5 and would put the same entry first about one time in four. Our three generations of one system passed 7, 4 and 8 of 45, a spread consistent with binomial noise.
The rules say the bundle checks that the corpus digest matches the release and that there are 3 attempts per task. In the commit we ran, it records the digest but does not compare it and does not check task or attempt counts; a 1-task bundle on a different corpus version reported no problems. The rules also describe a trace review by a reviewer, which we had no access to.
